% ============================================================
%  ch04  第一讲 §1.6 复数域上的仿射 Grassmann 簇、§1.7 $p$-进群
%  原文行 1192–1254
% ============================================================

\subsection{复数域上的仿射 Grassmann 簇}\label{top loop group}

本小节讨论复数域上的仿射 Grassmann 簇，故设 $k=\bC$。设 $\underline G=G\otimes_k \mO$，并把 $\Gr_{\underline G}$ 简记作 $\Gr$。此时 $\Gr$ 可以看作无穷维复解析空间（特别地，是 CW 复形）。

设 $K\subset G$ 是极大紧子群，$S^1$ 是单位圆。以 $\Omega K$ 记多项式映射 $(S^1,1)\to (K,e_K)$ 构成的空间。确切地说，用 $e^{i\phi}$（$\phi\in\bR$）参数化 $S^1$，并取定嵌入 $K\subset \on{SO}(n,\bR)$。则 $\Omega K$ 是这样的映射 $(S^1,1)\to (K,e_K)$ 构成的空间：与 $K\subset \on{SO}(n,\bR)$ 复合之后，映射的矩阵元都是 $e^{i\phi}$ 的 Laurent 多项式。把 $\Omega K$ 看作拓扑群。有下述定理。

\begin{thm}\label{top triv}~
\begin{enumerate}
\item 典范映射 $\Omega K\to LG$ 诱导同胚 $\Omega K\cong \Gr$。等价地说，有分解 $LG=\Omega K\cdot L^+G$。

\item 有同构 $\Omega K\times\Omega K\cong \Gr\tilde\times\Gr$，使下图交换
\[\begin{CD}
\Omega K\times \Omega K@>m>>\Omega K\\
@VVV@VVV\\
\Gr\tilde\times\Gr @>m>>\Gr,
\end{CD}\]
其中上行为环路群 $\Omega K$ 的通常乘法，下行为卷积映射 \eqref{conv m0}。
\end{enumerate}
\end{thm}

(1) 属于 Pressley--Segal（见 \cite[\S\ 8.3]{PS}）。严格说来，\emph{同上}只证明了 $G=\GL_n$ 的情形。但一般情形可由 Nadler 一个远为一般的定理推出（见 \cite[\S\ 4]{Na}）。确切地说，$\Omega K$ 可看作无穷维复群 $G(\bC[t,t^{-1}])$ 的一个「实形式」，Nadler 给出了 $G(\bC[t,t^{-1}])$ 的某些实形式在 $\Gr$ 上的轨道参数化；特殊到 $\Omega K$ 时，他的定理给出 (1)。(2) 由 (1) 直接推出。\zy{这里与拓扑环路群正式接上：nLab 词条 loop group（\url{https://ncatlab.org/nlab/show/loop+group}，2026-03-21 修订）记拓扑环路群的同胚分解 $\mathcal{L}G\simeq G\times\Omega G$；本文所得是更强的陈述 $\Gr\cong\Omega K$（A 级：\cite[\S 8.3]{PS}）。[来源：nLab loop group，2026-03-21；A 级 \cite{PS}。]}

\begin{rmk}
特别地，$\Gr$ 具有 $H$-空间的同伦型。此外，$L^+G$-主丛 $LG\to \Gr$ 平凡，且作为拓扑空间 $\Gr\tilde\times\Gr\cong \Gr\times\Gr$（不要把这个同构与 \eqref{I: isom} 混淆）。\S\ref{Lecture V} 将说明：对任意域 $k$ 上 $G$ 的仿射 Grassmann 簇，这些事实在上同调层面（事实上在 \'etale 同伦型层面）都有对应的体现。
\end{rmk}

\subsection{$p$-进群的仿射 Grassmann 簇}\label{p-adic}

简要讨论 $p$-进域 $F$（即 $\bQ_p$ 的有限扩张）上群的仿射 Grassmann 簇。与 $F=k((t))$ 的情形一样，第一步是让「由 $\bF_p$-代数 $R$ 的 $\Spec R$ 参数化的 $V=\bQ_p^n$ 的一族格」有意义；尤其要让圆盘 $D_R$ 有意义。虽然纤维积 $\Spec \bZ_p``\times''\Spec R$ 并无字面意义，但众所周知 $D_R$ 应当定义为
\[D_R:=\Spec W(R),\quad D^*_R:=\Spec W(R)[1/p],\]
其中 $W(R)$ 是 $R$ 的（$p$-典型）Witt 向量环。它的元素是序列 $(r_0,r_1,r_2,\ldots)\in R^{\bN}$，加法与乘法由某些（复杂的）多项式给出。到第一个分量的投影 $W(R)\to R$，$(r_0,r_1,\ldots)\mapsto r_0$，是满环同态，并带有乘法（但\emph{非加法}）截面 $R\to W(R)$，$r\mapsto [r]=(r,0,0,\ldots)$，称为 $r$ 的 \textbf{Teichm\"uller 提升}。若 $R$ 是完全环，$W(R)$ 中每个元素可唯一写成 $\sum_{i\geq 0} [a_i]p^i$，于是 $W(R)$ 可看作「以 $p$ 为变量、以 $R$ 为系数的幂级数环」。例如 $W(\bF_p)=\bZ_p$。一般介绍见 \cite[\S\ II.5]{Se2}。

一般地，设 $F$ 是 $p$-进域，$\mO$ 是其整数环，$k$ 是其剩余域。则可用 Drinfeld 引入的所谓分歧 Witt 向量类似地定义 $D_{F,R}$ 与 $D^*_{F,R}$。若 $R$ 完全，它们简单地由
\[D_{F,R}= \Spec W(R)\otimes_{W(k)}\mO,\quad D^*_{F,R}= \Spec W(R)\otimes_{W(k)}F\]
给出，其中 $W(k)\to \mO$ 是诱导剩余域恒等映射的唯一环同态。

现在设 $\underline G$ 是 $\mO$ 上光滑仿射群概形。则可完全仿照 \eqref{loc aff} 在 $\on{Aff}_k$ 上定义预层，但有一处需要注意：当 $R$ 不完全时，环 $W(R)$ 病态，例如 $R$ 不既约时 $p$ 可能是 $W(R)$ 的零因子。另一方面，注意

\begin{enumerate}
\item 若 $R$ 是完全 $k$-代数，则 $W(R)$ 性质良好，例如 $W(R)$ 中每个元素都有「幂级数展开」。
\item 概形 $X$ 在完全环上的 $R$-点决定了 $X$ 的完全化\footnote{完全 $k$-概形范畴是完全 $k$-代数上预层范畴的满子范畴，见 \cite[Lemma A.10]{Z14}。}。
\item 概形的（\'etale）拓扑（例如 $\ell$-进上同调）在过渡到完全化时不改变。
\end{enumerate}

因此，我们把 \eqref{loc aff} 中朴素的模空间问题限制到完全 $k$-代数范畴上。这样在该范畴上定义了预层；\'etale 拓扑与 fpqc 拓扑在该范畴上有意义，也可以在此框架下做代数几何（介绍见 \cite[Appendix A]{Z14}）。此时能问的最好问题是：这个函子是否可由完全 $k$-概形的归纳极限表示？下述困难定理（在 \cite{BS} 中证明）给出肯定回答。

\begin{thm}
函子 $\Gr_{\underline G}$ 可由归纳完全射影概形表示，称为 \textbf{Witt 向量仿射 Grassmann 簇}。
\end{thm}

想法是在 $\Gr_{\GL_n}$ 上构造丰沛的「行列式」线丛。困难在于：对 $R$-族格 $\La\subset W(R)^n$，商 $W(R)^n/\La$ \emph{不是} $R$-模，故 \S\ref{determinant line} 的构造不适用。此外，$\Gr_{\underline G}$ 上没有类似的 $\Aut(D)$ 对称性，$p$-进情形当然也没有定理 \ref{glob aff} 的类比。

此前在 \cite[\S~1]{Z14} 中证明了 $\Gr_{\underline G}$ 可由归纳完全真代数空间表示。那是个容易得多的定理，对大多数拓扑问题的应用（例如 \S\ref{Lecture V} 将讨论的几何 Satake 等价）已经够用。
